\section*{Chapitre \ref{ch:fm:the-risk-management-function} --- La fonction de gestion des risques}

\begin{solution}{exo:fm:the-risk-management-function:1}
Gravité 1 (le responsable de desk, le jour même)~; gravité 2 (le \omterm{def:fm:the-risk-management-function:cro}{directeur des risques}, dans un délai d'un jour ouvré)~; gravité 3 (le \omterm{def:fm:the-risk-management-function:cro}{comité des risques}, dans
un délai de deux jours ouvrés).
\end{solution}

\begin{solution}{exo:fm:the-risk-management-function:2}
De 2.3 à 4.0~milliards de dollars~: 74~\%. La limite de 2008 était de 1.5~milliards de dollars, soit 60~\% de plus que les 2.5~milliards de dollars de la méthode de l'année précédente.
\end{solution}

\begin{solution}{exo:fm:the-risk-management-function:3}
$(5+1)/2=3$ jours ouvrés~; $(10+1)/2=5.5$ avec une réunion tous les dix jours.
\end{solution}

\begin{solution}{exo:fm:the-risk-management-function:4}
Elle réécrit l'historique~: les dépassements survenus entre la date d'effet et la date d'approbation disparaissent d'un registre qui utilise la limite en vigueur, si bien
que personne ne peut voir plus tard que la société dépassait sa limite ou que la hausse a été accordée pour couvrir un dépassement. Un registre doit conserver les limites telles qu'elles étaient connues
chaque jour, avec les dates d'approbation et d'effet.
\end{solution}

\begin{solution}{exo:fm:the-risk-management-function:5}
Les desks qui demandent une hausse attendent la réunion suivante, en moyenne trois jours ouvrés, le dépassement restant ouvert. Lorsqu'ils doivent réduire pendant
l'attente, la plupart des dépassements se résorbent avant la réunion et la demande devient caduque~: la durée moyenne passe de 1.64 à 1.25 jours.
\end{solution}

\begin{solution}{exo:fm:the-risk-management-function:6}
Deux personnes interviennent, mais l'approbateur n'est pas indépendant~: le responsable de desk est rémunéré sur les résultats du desk et veut le même résultat que le
trader. Le second regard doit appartenir à quelqu'un dont les intérêts diffèrent, comme la fonction risques ou le \omterm{def:fm:the-risk-management-function:cro}{comité des risques}.
\end{solution}

\begin{solution}{exo:fm:the-risk-management-function:7}
Cinq dépassements, à $-51$, $-27$, $-9$, $-5$ et $-2$ jours ouvrés à compter de la modification. Celui à $-27$ est signalé pour une limite relevée pendant un
dépassement~; celui à $-5$, au moment où la première hausse a expiré, pour une limite relevée et des hausses répétées~; celui à $-2$, clos par le changement de modèle, pour
un modèle modifié pendant un dépassement et des hausses répétées. Les deux autres se sont refermés par des réductions et ne sont pas signalés.
\end{solution}

\begin{solution}{exo:fm:the-risk-management-function:8}
Le responsable du trading est le responsable des activités dont les risques sont mesurés~: l'objection passe par la personne qu'elle conteste, dont la rémunération
dépend des revenus que les limites contraignent. L'indépendance suppose une ligne hiérarchique vers le \omterm{def:fm:the-risk-management-function:cro}{comité des risques} et le directeur général, et non
une expertise du supérieur hiérarchique.
\end{solution}

\begin{solution}{pb:fm:the-risk-management-function:1}
\begin{enumerate}
\item Voir la \cref{def:fm:the-risk-management-function:cro}.
\item Une ligne hiérarchique extérieure aux activités, une rémunération indépendante de leurs revenus et un accès direct au \omterm{def:fm:the-risk-management-function:cro}{comité des risques}.
\item Indépendance et autorité, combinées à des échanges directs avec les responsables des activités, une escalade précoce (certaines dès l'été 2006) et
une information non filtrée par la hiérarchie.
\item Voir la \cref{def:fm:the-risk-management-function:foureyes}~; hausses de limites et modifications des modèles de risque.
\item 2.3~milliards de dollars~; 3.3~milliards de dollars pour l'exercice 2007 (fin 2006)~; 3.5~milliards de dollars le 7 septembre 2007~; 4.0~milliards de dollars approuvés le 14 janvier 2008,
avec effet au 3 décembre 2007.
\item Elle a retiré de l'historique les dépassements de la limite de 3.5~milliards de dollars entre le 3 décembre et l'approbation.
\item L'immobilier commercial, le capital-investissement et, un temps, les prêts à effet de levier~; une position de bridge en capitaux propres de 2.3~milliards de dollars, de mai à août 2007.
\item Preuves insuffisantes que la conduite sortait du champ de la business judgement rule ou relevait de la témérité~: les décisions imprudentes restaient du ressort de la direction.
L'autorité de la fonction risques doit être inscrite dans la procédure au préalable.
\item Voir la \cref{def:fm:the-risk-management-function:escalation}.
\item Gravité 1 jusqu'à 5~\% et deux jours, gravité 3 au-delà de 20~\% ou de cinq jours, gravité 2 entre les deux~; par exemple 3~\%/1 jour, 10~\%/1 jour et 3~\%/6 jours donnent 1, 2 et 3.
\item Voir la \cref{prop:fm:the-risk-management-function:wait}.
\item 45.3~\%, 14.4~\% et 5.4~\%.
\item 1.15, 1.64 et 1.25 jours~; 31.8, 43.5 et 34.6 jours de dépassement par an~; 2.44~\%, 1.33~\% et 0.74~\% des jours-desk.
\item Limites relevées ou modèles modifiés pendant un dépassement, et hausses répétées~; 54.7~\%, 23.1~\% et 11.9~\%.
\item La plupart des dépassements sont des dérives d'un jour qui s'inversent quel que soit le régime~; la durée moyenne, les jours de dépassement et les décomptes par gravité montrent la
différence.
\item Une hausse temporaire qui expire et est renouvelée aussitôt, puis un dépassement de 7.7~\% clos par un changement de modèle qui diminue de
un cinquième le risque mesuré sans aucune modification de la position.
\item Limites et modèles tels qu'ils étaient connus chaque jour, avec les dates d'approbation et d'effet pour chaque modification, et chaque signal d'alerte, qui ne doit jamais être effacé.
\item Voir la \cref{def:fm:the-risk-management-function:culture}~; la part des dépassements clos par des hausses, la part signalée, les dérogations au-delà de leur
échéance.
\item Responsable de desk approbateur~: 45.3~\% des dépassements clos par la hausse de la limite, médiane d'un jour (moyenne 1.15). Comité approbateur~: 14.4~\%, médiane égale à un
jour, mais une moyenne de 1.64 et davantage de jours de dépassement~; avec des réductions pendant l'attente, 5.4~\% et une moyenne de 1.25.
\item Les hausses et les changements de modèle exigent l'approbation du \omterm{def:fm:the-risk-management-function:cro}{comité des risques}, le desk réduisant vers sa limite pendant l'attente~; chaque changement est
enregistré avec sa date d'approbation, et les hausses répétées ainsi que les changements pendant un dépassement remontent au comité comme signaux d'alerte.
\end{enumerate}
\end{solution}

\begin{solution}{iq:fm:the-risk-management-function:1}
Au \omterm{def:fm:the-risk-management-function:cro}{comité des risques} du conseil d'administration et au directeur général, et non au responsable des activités dont il mesure le risque, afin que les objections ne puissent
pas être filtrées ni écartées par ceux qu'elles concernent.
\iqlookfor{indépendance de la ligne hiérarchique et accès direct au conseil d'administration.}
\end{solution}

\begin{solution}{iq:fm:the-risk-management-function:2}
Il est consigné et classé~; le responsable de desk en est informé le jour même~; le desk réduit ou demande une hausse temporaire approuvée de façon indépendante, avec une
échéance~; s'il persiste, le dépassement est remonté au \omterm{def:fm:the-risk-management-function:cro}{directeur des risques}.
\iqlookfor{la procédure, et pas seulement la réduction.}
\end{solution}

\begin{solution}{iq:fm:the-risk-management-function:3}
Refuser par défaut le renouvellement et le porter au comité~: soit la limite est erronée et doit être réinitialisée par le processus annuel, soit le desk
utilise les hausses pour porter plus de risque que sa limite ne le prévoyait.
\iqlookfor{hausses répétées comme signal d'alerte.}
\end{solution}

\begin{solution}{iq:fm:the-risk-management-function:4}
Le dépassement subsiste sous l'ancien modèle jusqu'à la validation indépendante du nouveau~; la validation examine le modèle sur le fond, et non sur son effet,
et le changement est enregistré avec sa date et signalé.
\iqlookfor{validation indépendante du dépassement.}
\end{solution}

\begin{solution}{iq:fm:the-risk-management-function:5}
La part des dépassements clos par des hausses, les changements effectués pendant les dépassements, les changements antidatés, les dérogations au-delà de leur échéance et le temps que mettent les mauvaises
nouvelles à parvenir au comité.
\iqlookfor{des tendances, non des décomptes.}
\end{solution}

\begin{solution}{iq:fm:the-risk-management-function:6}
Tables en ajout seul des limites, des versions de modèles et des expositions, chacune avec des horodatages d'effet et d'enregistrement~; dépassements dérivés de ces tables à la
date de chaque jour~; approbations conservées comme enregistrements signés distincts~; signaux d'alerte stockés et jamais supprimés.
\iqlookfor{enregistrements bitemporels et stockage en ajout seul.}
\end{solution}
